\documentclass[a4paper,12pt]{article}
\usepackage{myStyle}

%%%%%%%%%%%%%%%%%%%%%%%%%%%%%%%%%%%%%%%%%%%%%%%%%%%%%%%%%%%%%%%%%%%%%
% Hier eigene Daten einfügen                                        %
%%%%%%%%%%%%%%%%%%%%%%%%%%%%%%%%%%%%%%%%%%%%%%%%%%%%%%%%%%%%%%%%%%%%%
\newcommand{\Jahr}{2012 / 2013}         % Typ: "2011 / 2012" oder "2012"
\newcommand{\Semester}{Wintersemester}  % "Wintersemester" oder "Sommersemester"
\newcommand{\Datum}{\today}             % Wann wurde der Bericht erstellt?
\newcommand{\Ort}{Karlsruhe}
\newcommand{\Vorname}{Max}
\newcommand{\Nachname}{Mustermann}
\newcommand{\Matrikelnummer}{123456}
\newcommand{\Studienfach}{Informatik (BA)}
\newcommand{\Institut}{Institut für Theoretische Informatik}
\newcommand{\TutoriumTitle}{Programmieren Tutorium}
\newcommand{\Email}{info@martin-thoma.de}
%%%%%%%%%%%%%%%%%%%%%%%%%%%%%%%%%%%%%%%%%%%%%%%%%%%%%%%%%%%%%%%%%%%%%
\hypersetup{ 
  pdfauthor   = {\Vorname~\Nachname}, 
  pdfkeywords = {KIT; \Vorname~\Nachname}, 
  pdftitle    = {Abschlussreflexion von~\Vorname~\Nachname~-~\Semester~\Jahr} 
} 
 
\pagestyle{fancy}
\fancyhf{}
\renewcommand{\headrulewidth}{0pt}
\renewcommand{\footrulewidth}{0pt}
\fancyfoot[R]{Seite~\thepage~von \pageref{LastPage}}
 
\definecolor{LightCyan}{rgb}{0.88,1,1}
\definecolor{LightGreen}{rgb}{0.68,1,0.68}
\definecolor{Rosa}{rgb}{1,0.68,0.68}
 
\pagenumbering{arabic}
 
\begin{document}
 \pagenumbering{roman}
 \begin{titlepage}
    \author{Martin Thoma}
	\ifAFive
    \title{Geometrie und Topologie\\\vspace{4cm}
        \includegraphics[width=0.9\linewidth]{figures/Torus.pdf}}
	\else
    \title{Geometrie und Topologie\\\vspace{10cm}
        \includegraphics[width=0.9\linewidth]{figures/Torus.pdf}}	
	\fi
    \maketitle
\end{titlepage}

 \setcounter{page}{2}
 \pagenumbering{arabic}

% \title{Semesterbericht über das \Semester \Jahr}
% \author{\Vorname \Nachname}
% \date{\Datum}

\section*{Meine Lehre und ihr Umfeld}
\subsection*{Rahmenbedingungen}
Das Programmieren-Tutorium habe ich jeden Montag von 15:45 bis 17:15 
Uhr gehalten. Insgesamt habe ich so an 14 Tagen meinen 19 Studenten 
Grundlagen in Java beigebracht. Allerdings waren von den 19 Studenten
nur etwa sieben regelmäßig anwesend.\\

\clearpage
\section*{Ablaufplanung für eine Einzelveranstaltung}
\textbf{Datum und Thema}: TODO\\
\textbf{Lernziele}: TODO\\

\begin{tabularx}{\textwidth}{@{}llllX}
Dauer   & Inhalte   & Methoden  & Medien    & Sinn \& Zweck\\
\hline
1 min   & Begrüßung &   -       & Folien    & -\\
1 min   & Begrüßung &   -       & Folien    & -\\
1 min   & Begrüßung &   -       & Folien    & -\\
1 min   & Begrüßung &   -       & Folien    & -\\
1 min   & Begrüßung &   -       & Folien    & -\\
\end{tabularx}

\textbf{Kurzreflexion}:

(mind. 0.5 Seiten, siehe Fragen unter 4.)

\noindent \Ort, den \Datum\\
\\
\includegraphics[height=10mm]{max-mustermann.pdf}\\
\Vorname~\Nachname
\end{document}
